%% kapa: exam=1 round=31 year=2020 period=1 form=A subject=real_estate subject_ko=부동산학원론 seq=36 number=116
%% source: Q-Net 감정평가사 기출문제 (cst00309) artlSeq=5207219 file=extracted/1008/5207219/A형/A형/2020년도 제31회 감정평가사 1차 시험_1교시_A형.pdf
%% conversion: kordoc 4.24.51 -> tools/gen1.py (auto); review=proofread-sonnet; flags=none
\begin{question}{116}
부동산가치의 발생요인에 관한 설명으로 옳지 않은 것은?
\begin{choices}
  \item 유효수요는 구입의사와 지불능력을 가지고 있는 수요이다.
  \item 효용(유용성)은 인간의 필요나 욕구를 만족시킬 수 있는 재화의 능력이다.
  \item 효용(유용성)은 부동산의 용도에 따라 주거지는 쾌적성, 상업지는 수익성, 공업지는 생산성으로 표현할 수 있다.
  \item 부동산은 용도적 관점에서 대체성이 인정되고 있기 때문에 절대적 희소성이 아닌 상대적 희소성을 가지고 있다.
  \item 이전성은 법률적인 측면이 아닌 경제적인 측면에서의 가치발생요인이다.
\end{choices}
\end{question}
\answer{5}
